\vspace{-8pt}
\section{Experiments}
\label{sec:exp}
\vspace{-10pt}
We aim to answer the following questions: \textbf{RQ1 (Effectiveness)}: How does \hero’s overall performance compare to state-of-the-art (SOTA) baselines? (\S \ref{sec:overall_perf}) \textbf{RQ2 (Component Contribution)}:  What is the individual contribution of the key components? (\S\ref{sec:ablation}) \textbf{RQ3 (Efficiency)}: Does improved performance, if any, come at the cost of increased token consumption?(\S\ref{sec:token_usage}) \textbf{RQ4 (Topology Evolution)}: How do agent coordination topologies evolve as the system accumulates experience and updates agent prompts? (\S\ref{sec:topology_analysis})



\textbf{Tasks and Datasets.}We evaluate \hero on multi-hop QA benchmarks 2WikiMultiHopQA (2WikiQA) \citep{ho-etal-2020-constructing}, HotpotQA \citep{yang2018hotpotqa}, MusiQue \citep{trivedi2021musique} and AmbigQA \citep{min-etal-2020-ambigqa} for ambiguous QA. For out-of-distribution (OOD) evaluation, we use Bamboogle \citep{press-etal-2023-measuring} and HoVer \citep{jiang2020hover} for multi-hop QA and fact verification respectively. Evaluation metrics involve Exact Match (EM), accuracy, and F1 score for QA tasks, and accuracy for fact verification.

\textbf{Baselines.}
We compare \hero against:
(1) \textbf{Direct inference and chain-of-thought (CoT) without RAG} across different LLM families and sizes;
(2) \textbf{Single-turn RAG} with different LLMs as answer generators; (3) \textbf{Iterative RAG}: IterDRAG \citep{yue2024inference}, Plan-RAG \citep{verma2024plan}
Search-o1 \citep{li2025search}, Search-R1 \citep{jin2025search}, IRCoT \citep{trivedi-etal-2023-interleaving}, SELF-RAG \citep{asai2023self}, CORAG \citep{wang2025chain}; and (4) \textbf{Agentic RAG}, including InstructRAG, R1-Searcher \citep{song2025r1old},  DeepResearcher \citep{zheng-etal-2025-deepresearcher}, MMOA-RAG \citep{chen2025mao}, AceSearcher \citep{xu2025acesearcher} , and MAO-ARAG \citep{chen2025mao}, ExSearch \citep{shi2025iterative}.
Model and implementation details are provided in Appendix~\ref{app:baselin_model}.

%, including Qwen-3-4B, 8B, and 14B \citep{yang2025qwen3}, Qwen2.5-Instruct-7B \citep{hui2024qwen2}, LLaMA3.1-8B \citep{grattafiori2024llama}, and GPT-4o-mini \citep{hurst2024gpt};


\textbf{Implementation Details.} We consider core agents commonly employed in RAG: Query Decomposer, Retriever, Answer Generator, Query Rewriter, Evidence Selector, Context Validator, Reflect Agent, and Conclude Agent. We use Wikipedia as the corpus, returning top-\emph{5} passages per step for multi-step retrieval and top-\emph{10} for single-turn methods via the BGE retriever \citep{xiao2024c}. We use \llama \citep{grattafiori2024llama} and \qwen \citep{yang2025qwen3} as the backbone LLMs for the orchestrator and GPT-4o-mini for agents. All backbones remain frozen during learning and inference. To highlight data efficiency, we first use GPT-4o to annotate queries from 2WikiQA, HotpotQA, AmbigQA, and MusiQue along \textbf{reasoning type} (bridge, intersection, comparison, temporal multi-hop, ambiguous) and \textbf{complexity} (easy, medium, hard). Based on these annotations, we perform stratified, difficulty-aware sampling to construct a curated training set that ensures balanced coverage across reasoning categories while preserving distributional diversity (statistics in Appendix~\ref{app:training_set}). Generative sampling uses temperature 0.9 during group rollouts, and $0.0 / 0.3$ for in-distribution and OOD evaluation, respectively.  (Pseudocode for \hero is provided in Appendix~\ref{app:pseducode}).
